\documentclass[../../../include/open-logic-section]{subfiles}

\begin{document}

\olfileid{his}{set}{pathology}
\olsection{অস্বাভাবিক নির্মাণ}

কিন্তু \emph{সীমা}-র সংজ্ঞা থেকে
কিছু বেশ “অস্বাভাবিক” নির্মাণের
সম্ভাবনাও দেখা গেল।

১৮৩০-এর দশক নাগাদ বোল্‌ৎসানো
একটি অপেক্ষক খুঁজে পান যা
\emph{সর্বত্র অবিচ্ছিন্ন}, কিন্তু
\emph{কোথাও অবকলনযোগ্য নয়}।
(দুর্ভাগ্যবশত বোল্‌ৎসানো
এটি প্রকাশ করেননি; একই ধারণা
ভাইয়ারস্ট্রাস স্বাধীনভাবে আবিষ্কার
করার ফলে গণিতবিদেরা প্রথম
১৮৭২ সালে এর পরিচয় পান।)\footnote{এই
ইতিহাস \href{http://en.wikipedia.org/wiki/Weierstrass_function}{ভাইয়ারস্ট্রাস
অপেক্ষক} নিয়ে উইকিপিডিয়ার
নিবন্ধের অত্যন্ত বিস্তারিত
পাদটীকায় নথিবদ্ধ আছে।}
ঘটনাটি অন্তত বিস্ময়কর ছিল।
$|x|$-এর মতো সর্বত্র
অবিচ্ছিন্ন অথচ কোনো
একটি নির্দিষ্ট বিন্দুতে
অবকলনযোগ্য নয়—এমন
অপেক্ষক সহজেই পাওয়া যায়।
কিন্তু সর্বত্র অবিচ্ছিন্ন
অথচ \emph{কোথাও}
অবকলনযোগ্য নয়—এমন
অপেক্ষক একেবারেই
অন্য রকম। একবার
ভেবে দেখো, এমন
অপেক্ষকের রেখাচিত্র
কীভাবে আঁকতে
চাইবে। অবিচ্ছিন্ন
হতে হলে কাগজ
থেকে কলম না তুলেই
আঁকা সম্ভব হতে
হবে; কিন্তু কোথাও
অবকলনযোগ্য না
হতে হলে কলমের
দিক নিরন্তর
হঠাৎ বদলাতে হবে।

এর পরেও আরও
“অস্বাভাবিক” ফল
এল। ১৮৭৪ সালের
৫ জানুয়ারি কান্টর
ডেডেকিন্ডকে একটি
চিঠিতে প্রশ্ন
করেন:
\begin{quote}
একটি তলের (যেমন
সীমানাসহ বর্গক্ষেত্র)
প্রতিটি বিন্দুর সঙ্গে
একটি রেখার (যেমন
দুই প্রান্তবিন্দুসহ
একটি সরলরেখা)
একটি করে বিন্দু
এমনভাবে একৈক
সম্পর্কে আনা
যায় কি, যাতে
তলের প্রতিটি
বিন্দুর জন্য
রেখায় একটি
বিন্দু থাকে,
এবং বিপরীতভাবে
রেখার প্রতিটি
বিন্দুর জন্য
তলে একটি
বিন্দু থাকে?

এখনও আমার মনে
হয়, এই প্রশ্নের
উত্তর দেওয়া
খুব কঠিন—যদিও
এ ক্ষেত্রেও
\emph{না} বলার
তাগিদ এত
প্রবল যে
প্রমাণটিকে
প্রায় অনাবশ্যক
মনে হয়।
[\citealt{Gouvea2011}-এ
উদ্ধৃত]
\end{quote}
কিন্তু ১৮৭৭ সালে
কান্টর প্রমাণ
করেন, তিনি
ভুল ছিলেন।
বাস্তবে একটি
রেখা ও একটি
বর্গক্ষেত্রে
বিন্দুর সংখ্যা
ঠিক সমান।
১৮৭৭ সালের
২৯ জুন
ডেডেকিন্ডকে
তিনি লেখেন,
“\emph{je le vois,
mais je ne le
crois pas}”;
অর্থাৎ, “আমি
এটি দেখছি,
কিন্তু বিশ্বাস
করতে পারছি
না।” গণিতের
“প্রচলিত ইতিহাসে”
এটি প্রায়ই
তৎকালীন
গণিতবিদদের
কাছে নতুন
ফলগুলি কতটা
\emph{আক্ষরিক
অর্থেই অবিশ্বাস্য}
ছিল তার
ইঙ্গিত
হিসেবে ধরা
হয়। (চিঠিপত্র
\citet{Gouvea2011}-এ
দেওয়া আছে;
আমরা
\olref[his][set][mythology]{sec}-তে
ফিরে আসব।
কান্টরের
প্রমাণের
রূপরেখা আছে
\olref[his][set][cantorplane]{sec}-তে।)

কান্টরের ফল
থেকে অনুপ্রাণিত
হয়ে পেয়ানো
ভাবতে শুরু
করেন, একটি
রেখাকে
\emph{অবিচ্ছিন্নভাবে}
পুরো তলের
উপর চিত্রিত
করা যায়
কি না। তা
হলে সেটি
\emph{স্থানপূরক
বক্ররেখা}
হবে।
\citeyear{Peano1890}
সালে পেয়ানো
ঠিক এমনই
একটি বক্ররেখা
নির্মাণ করেন।
% BN-SRC-858 TARGET-CORRECTION-BEGIN
\par\noindent\textbf{পাঠ-সংশোধন।} মূলপাঠের “smoothly” এখানে অবিচ্ছিন্ন চিত্রণের অর্থে
অনুবাদ করা হয়েছে। পেয়ানোর মূল প্রবন্ধে অপেক্ষকগুলি
অবিচ্ছিন্ন এবং কোথাও অবকলনযোগ্য নয়; দেখো
\href{https://webhomes.maths.ed.ac.uk/~v1ranick/papers/peano.pdf}{মূল প্রবন্ধ}, পৃ.~১৫৭ ও ১৬০।
% BN-SRC-858 TARGET-CORRECTION-END
এটি সত্যিই
স্বজ্ঞাবিরোধী:
ইউক্লিড রেখাকে
বলেছিলেন
“প্রস্থহীন
দৈর্ঘ্য”
(প্রথম খণ্ড,
সংজ্ঞা ২);
কিন্তু পেয়ানো
দেখান যে
রেখাকে
যথাযথভাবে
পেঁচিয়ে
দৈর্ঘ্যকে
প্রস্থে
রূপান্তর
করা যায়।
\citeyear{Hilbert1891}
সালে হিলবার্ট
তুলনায় কিছুটা
সহজবোধ্য
একটি
স্থানপূরক
বক্ররেখার
বর্ণনা দেন,
সঙ্গে তার
কয়েকটি
ছবিও।
বক্ররেখাটি
ধাপে ধাপে
নির্মিত;
নিচে
নির্মাণের
প্রথম
ছয়টি
ধাপ:
\begin{center}
\begin{tikzpicture}
\begin{scope}[xshift=20pt, yshift=20pt]
	\draw [oldiagcolorC, l-system={Hilbert curve, step=40pt, angle=90, axiom=L, order=1}] lindenmayer system;
\end{scope}
\begin{scope}[xshift=110pt, yshift=10pt]
	\draw [oldiagcolorC, l-system={Hilbert curve, step=20pt, angle=90, axiom=L, order=2}] lindenmayer system;
\end{scope}
\begin{scope}[xshift=205pt, yshift=5pt]
	\draw [oldiagcolorC, l-system={Hilbert curve, step=10pt, angle=90, axiom=L, order=3}] lindenmayer system;
\end{scope}
\begin{scope}[xshift=2.5pt, yshift=-97.5pt]
	\draw [oldiagcolorC, l-system={Hilbert curve, step=5pt, angle=90, axiom=L, order=4}] lindenmayer system;
\end{scope}
\begin{scope}[xshift=101.25pt, yshift=-98.75pt]
	\draw [oldiagcolorC, l-system={Hilbert curve, step=2.5pt, angle=90, axiom=L, order=5}] lindenmayer system;
\end{scope}
\begin{scope}[xshift=200.625pt, yshift=-99.375pt]
	\draw [oldiagcolorC, l-system={Hilbert curve, step=1.25pt, angle=90, axiom=L, order=6}] lindenmayer system;
\end{scope}
\draw[gray] (0pt,0pt) rectangle (80pt, 80pt);
\draw[gray] (100pt,0pt) rectangle (180pt, 80pt);
\draw[gray] (200pt,0pt) rectangle (280pt, 80pt);
\draw[gray] (0pt,-100pt) rectangle (80pt, -20pt);
\draw[gray] (100pt,-100pt) rectangle (180pt, -20pt);
\draw[gray] (200pt,-100pt) rectangle (280pt, -20pt);
\end{tikzpicture}
\end{center}
সীমায়—যে
ধারণার
কঠোর
সংজ্ঞা
ততদিনে
পাওয়া
গিয়েছিল—পুরো
বর্গক্ষেত্র
!!{colorC}
রঙে
ভরে যায়।
প্রসঙ্গত,
হিলবার্টের
বক্ররেখাটি
সর্বত্র
অবিচ্ছিন্ন,
কিন্তু কোথাও
অবকলনযোগ্য
নয়; স্বজ্ঞায়
কারণ, অসীম
সীমায়
অপেক্ষকটি
প্রতিটি
মুহূর্তে
হঠাৎ দিক
বদলায়।
(হিলবার্টের
নির্মাণ
আরও বিশদে
দেখব
\olref[his][set][hilbertcurve]{sec}-তে।)

ভালোই হোক
বা মন্দই
হোক,
এমন
“অস্বাভাবিক”
জ্যামিতিক
নির্মাণকে
জ্যামিতিক
স্বজ্ঞার
উপর
নির্ভরতা
সন্দেহের
চোখে
দেখার
কারণ
হিসেবে
নেওয়া
হয়েছিল।
গণিতের
এক নতুন
পদ্ধতির
পক্ষে
এগুলি
প্রায়
\emph{প্রচারণার}
মতো হয়ে
ওঠে;
শেষে
যা
সেটতত্ত্বে
পরিণতি
পায়।
পরে
বিষয়টি
নিয়ে
পুরাণকথা
গড়তে
গিয়ে
বারবার
বলা
হয় যে
ফলগুলি
যেমন
সম্পূর্ণ
কঠোর,
তেমনই
সম্পূর্ণ
চমকপ্রদ।
সুতরাং
এগুলি
দুই
কাজে
লাগে:
জ্যামিতিক
স্বজ্ঞার
উপর
নির্ভর
না করার
সতর্কতা
হিসেবে,
এবং
নতুনভাবে
চিন্তার
উর্বরতা
দেখানোর
উদাহরণ
হিসেবে।

\end{document}
